% TDT4225 Assignment 2 - report
% Follows "Report - template - Assignment 2.docx":
%   A4, 1 inch margins, Lato, NTNU logo left + TDT4225 right in the header
%   with a grey rule below, page number centred in the footer,
%   title 18 pt bold, subtitle 14 pt, headings 14 pt bold, body 12 pt.
\documentclass[12pt,a4paper]{article}

\usepackage[T1]{fontenc}
\usepackage[utf8]{inputenc}
\usepackage[default]{lato}
\usepackage[a4paper,margin=1in,headheight=28pt,headsep=18pt]{geometry}
\usepackage{graphicx}
\usepackage{xcolor}
\usepackage{fancyhdr}
\usepackage{titlesec}
\usepackage{booktabs}
\usepackage{tabularx}
\usepackage{float}
\usepackage{listings}
\usepackage{enumitem}
\usepackage[hidelinks]{hyperref}

\setlength{\parindent}{0pt}
\setlength{\parskip}{6pt}
\setlist{nosep,leftmargin=*}

% ---- header / footer as in the template ----
\pagestyle{fancy}
\fancyhf{}
\fancyhead[L]{\raisebox{-0.2\height}{\includegraphics[width=3.86cm]{ntnu_logo.png}}}
\fancyhead[R]{TDT4225}
\renewcommand{\headrule}{{\color[HTML]{A0A0A0}\hrule height 1.5pt}}
\fancyfoot[C]{\thepage}

% ---- headings: plain bold 14 pt (Introduction, Results, ...) ----
\titleformat{\section}{\fontsize{14}{17}\selectfont\bfseries}{}{0pt}{}
\titleformat{\subsection}{\fontsize{12}{15}\selectfont\bfseries}{}{0pt}{}
\titlespacing*{\section}{0pt}{18pt}{6pt}
\titlespacing*{\subsection}{0pt}{12pt}{4pt}

\lstset{basicstyle=\ttfamily\footnotesize,columns=fullflexible,
        keepspaces=true,frame=single,rulecolor=\color[HTML]{A0A0A0},
        breaklines=true}

% Screenshot of the console output. Put the PNGs in screenshots/.
% A missing file shows a grey placeholder box instead of breaking the build.
\newcommand{\screenshot}[2]{%
  \begin{figure}[H]\centering
  \IfFileExists{screenshots/#1.png}%
    {\includegraphics[width=\linewidth,height=0.8\textheight,keepaspectratio]{screenshots/#1.png}}%
    {\fbox{\parbox[c][2.5cm][c]{0.9\linewidth}{\centering\color{gray}
      screenshot missing: screenshots/#1.png}}}%
  \caption{#2}\label{fig:#1}
  \end{figure}}

% EDA figure, optional width (fraction of \linewidth), label fig:<name>
\newcommand{\edafig}[3][0.8]{%
  \begin{figure}[H]\centering
  \includegraphics[width=#1\linewidth]{figures/#2.pdf}
  \caption{#3}\label{fig:#2}
  \end{figure}}

\begin{document}

% ---- title block ----
\vspace*{0pt}
\begin{center}
  {\fontsize{18}{22}\selectfont\bfseries Report\par}
  \vspace{4pt}
  {\fontsize{14}{17}\selectfont Assignment 2 -- MySQL\par}
\end{center}
\vspace{12pt}

\textbf{Group}: 16\\
\textbf{Students}: Ingrid Talseth Singstad, Olav Larsen Halleraker

\vspace{12pt}

% =====================================================================
\section*{Introduction}

In this assignment we worked with the Porto Taxi Trajectory dataset:
1.7 million taxi trips from July 2013 to June 2014, where each trip has a
GPS point every 15 seconds. We first explored and cleaned the data, then
designed a MySQL schema and loaded the cleaned data into it, and finally
answered ten questions with SQL and Python.

The main problems we had to solve were that TRIP\_ID is not unique, that
some trips have almost no GPS points or obviously wrong points, and that
83 million GPS points take a lot of space and time to load.

During the first days we sat together at school, went through the assignment and agreed on 
a shared plan for the EDA, the data cleaning, the schema and each of the 
queries. We wrote the first draft of the code together, using AI-assisted
development, and then divided the rest of the implementation between us.
Ingrid continued with the EDA, the data loading and the queries, and
Olav took over the final stage, completing the remaining parts and
getting everything ready for delivery. Whenever one of us had questions
about the other's work, we discussed it together, either online or in person.

MySQL 8.0.39 ran locally in Docker, and we used Python with
\texttt{mysql-connector-python}, \texttt{pandas} and \texttt{haversine}.

Repository: \url{https://github.com/olavhalleraker/tdt4225-project}

% =====================================================================
\section*{Results}

\subsection*{Part 1 -- Exploratory data analysis}

The raw CSV has \textbf{1,710,670 trips}, \textbf{83,409,386 GPS points}
and \textbf{448 taxis}, covering exactly one year.

\textbf{Metadata.}
\begin{itemize}
  \item ORIGIN\_CALL and ORIGIN\_STAND are the only columns with empty values
        (78.7\,\% and 52.9\,\%). This is expected, since they only apply to
        call type A and B. 11,302 B trips still have no stand.
  \item TRIP\_ID is TIMESTAMP, the digit 6 and TAXI\_ID put together. 80 IDs
        occur more than once (161 rows). In 35 of them a stub of 0--6 points
        sits next to a longer trip, and in 31 all copies have fewer than
        3 points.
  \item Call types: A 21.3\,\%, B 47.8\,\%, C 30.9\,\%.
  \item DAY\_TYPE is ``A'' in every row, even on Christmas Day, and
        MISSING\_DATA is true for only 10 rows. Neither column carries
        usable information.
  \item TIMESTAMP is Unix time, i.e.\ UTC. 157,929 trips end up in a
        different time band depending on whether we use UTC or Porto time,
        so we use local Porto time (Europe/Lisbon) everywhere.
  \item Trips start at all hours, but between 09 and 18 there are on average
        1.7 times as many per hour as between 00 and 08
        (Figure~\ref{fig:metadata_trips_by_hour}).
\end{itemize}

\textbf{Trajectories.}
\begin{itemize}
  \item Median 41 points per trip (Figure~\ref{fig:polyline_points_per_trip}).
        43,904 trips (2.6\,\%) have fewer than 3 points: 5,901 are empty,
        30,609 have one point and 7,394 have two.
  \item Median duration is 10 minutes and median distance 4.0\,km, but the
        longest trip is 16\,h and the longest distance 1,229\,km.
  \item 32 trips have points far outside Portugal, up to 51$^\circ$N and
        53$^\circ$E. These are considered GPS errors.
  \item 38,216 trips have at least one jump of more than 833\,m in 15\,s
        (above 200\,km/h, Figure~\ref{fig:polyline_max_speed}). Most of them
        look like a gap in the signal rather than a single bad point: we
        found only 208 single points that jump away and straight back. These 
        jumps explain the extreme distances.
  \item 7,142 trips leave the Porto area
        (Figure~\ref{fig:polyline_map_hexbin}). Apart from the 32 above,
        85\,\% of them have no jump at all, so they seem like real trips on
        the motorway, and we kept them.
\end{itemize}

\edafig{polyline_points_per_trip}{Number of GPS points per trip.}
\edafig{metadata_trips_by_hour}{Trips per hour of the day (local time).}
\edafig{polyline_max_speed}{Highest speed between two consecutive points in each valid trip.}
\edafig[0.6]{polyline_map_hexbin}{Density of GPS points around Porto. The red cross marks City Hall.}

\textbf{Missing and buggy data.} Table~\ref{tab:cleaning} sums up what we
found and how we handled it. We removed as little as possible and marked the
rest, so the queries can decide what to include.

\begin{table}[H]\centering\small
\renewcommand{\arraystretch}{1.2}
\begin{tabularx}{\linewidth}{>{\raggedright\arraybackslash}X r >{\raggedright\arraybackslash}X r}
\toprule
Problem & Rows & What we did & Removed\\
\midrule
Empty ORIGIN\_CALL / ORIGIN\_STAND (only used by call type A / B) & -- & Stored as NULL & 0\\
B trips without ORIGIN\_STAND & 11,302 & Kept, stand is NULL & 0\\
Duplicate TRIP\_ID (80 IDs) & 161 & Kept the copy with the most GPS points, or the first one in the CSV on a tie. A taxi cannot start two trips in the same second, so the other copies (0--6 points each) are recording errors, not real trips & 81\\
Fewer than 3 GPS points (5,901 empty) & 43,904 & Kept and marked \texttt{is\_valid = 0}, since Q7 has to count them & 0\\
GPS points outside mainland Portugal & 32 & Removed the whole trip. It is only 0.002\,\% of the trips, so the effect on the answers is negligible, and it is simpler than repairing them & 32\\
Jump above 200\,km/h between two points & 38,216 & Kept and marked \texttt{has\_gap = 1}, so distance queries can leave them out & 0\\
DAY\_TYPE is always ``A'' & all & Stored, but not used & 0\\
MISSING\_DATA is true for only 10 trips & 10 & Stored, but not used. We judge gaps from the trajectories instead & 0\\
Very long trips (more than 3\,h) & 847 & Kept. Some may be a meter left running, but we cannot tell which, and they are only 0.05\,\% of the trips & 0\\
\bottomrule
\end{tabularx}
\caption{Missing and buggy data. Rows are counted in the raw CSV.}
\label{tab:cleaning}
\end{table}

After cleaning the database has \textbf{448 taxis, 1,710,557 trips and
83,408,413 GPS points}. Only 113 rows were removed in total. The counts
of invalid trips and trips with a jump are a little lower after cleaning
(43,815 and 38,198), because some of the removed rows were among them.

\screenshot{eda_output_1}{EDA output (\texttt{eda/eda.py}), checks 3, 5 and 6.}
\screenshot{eda_output_2}{EDA output (\texttt{eda/eda.py}), checks 9, 12 and 13.}

\subsection*{Part 1 -- Schema}

We use four tables: \texttt{call\_type}, \texttt{taxi}, \texttt{trip} and
\texttt{gps\_point}. The most important choices were:
\begin{itemize}
  \item \texttt{trip} gets its own 4-byte integer key instead of the
        19-digit TRIP\_ID, which needs an 8-byte BIGINT. TRIP\_ID is unique
        after cleaning (we enforce it with a UNIQUE key), but the trip key is
        repeated in all 83 million rows of \texttt{gps\_point}, so the
        smaller key saves about 320\,MB there.
  \item Values we need often (start and end time, duration, distance,
        number of points, start and end point, bounding box) are computed
        once when loading. That way, only Q6 has to read the 83 million
        GPS points. These columns repeat data from \texttt{gps\_point}, but
        the data is loaded once and never updated, so they cannot get out
        of sync.
  \item \texttt{gps\_point} is kept as small as possible: a composite primary
        key, no timestamp per point (it is \texttt{start\_time + 15 s $\cdot$ seq})
        and no extra indexes. The whole database is 3.3\,GB.
  \item Coordinates are stored as DECIMAL(8,6) and DECIMAL(9,6). The CSV
        never has more than 6 decimals, so every value is stored exactly, in
        4 and 5 bytes. FLOAT also takes 4 bytes, but cannot store every
        6-decimal value exactly.
  \item Deleting a taxi deletes its trips, and deleting a trip deletes its
        GPS points (\texttt{ON DELETE CASCADE}), so a bad trip can be removed
        with one statement without leaving orphan rows. \texttt{call\_type}
        uses \texttt{RESTRICT}, so a code cannot be removed while trips use
        it.
\end{itemize}

The loader (\texttt{insert\_data.py}) reads the CSV twice. The first pass
decides which rows to keep, and the second parses the trips in batches and
bulk-loads them with \texttt{LOAD DATA LOCAL INFILE}, sent from Python
through \texttt{mysql-connector-python} instead of row-by-row
\texttt{INSERT}s as in \texttt{example.py}. The full load takes about
4 minutes. The foreign key and unique checks are turned off during the
load, so afterwards we checked that the number of points matches
\texttt{SUM(n\_points)} (shown in the Q1 output, Figure~\ref{fig:q1}), that
all trip codes are distinct, that there are no orphan rows, and that a few
trips match the CSV exactly.

\begin{lstlisting}[language=SQL]
CREATE TABLE call_type (
    code        CHAR(1)      NOT NULL,
    description VARCHAR(40)  NOT NULL,
    PRIMARY KEY (code)
);

CREATE TABLE taxi (
    taxi_id     INT UNSIGNED NOT NULL,
    PRIMARY KEY (taxi_id)
);

CREATE TABLE trip (
    id            INT UNSIGNED      NOT NULL AUTO_INCREMENT,
    trip_code     BIGINT UNSIGNED   NOT NULL,
    taxi_id       INT UNSIGNED      NOT NULL,
    call_type     CHAR(1)           NOT NULL,
    origin_call   INT UNSIGNED      NULL,
    origin_stand  SMALLINT UNSIGNED NULL,
    day_type      CHAR(1)           NOT NULL,
    missing_data  BOOLEAN           NOT NULL,
    start_unix    INT UNSIGNED      NOT NULL,
    start_time    DATETIME          NOT NULL,
    end_time      DATETIME          NOT NULL,
    duration_s    INT UNSIGNED      NOT NULL,
    n_points      SMALLINT UNSIGNED NOT NULL,
    is_valid      BOOLEAN AS (n_points >= 3) VIRTUAL,
    distance_km   DECIMAL(9,3)      NOT NULL,
    has_gap       BOOLEAN           NOT NULL DEFAULT 0,
    start_lat     DECIMAL(8,6)      NULL,
    start_lon     DECIMAL(9,6)      NULL,
    end_lat       DECIMAL(8,6)      NULL,
    end_lon       DECIMAL(9,6)      NULL,
    min_lat       DECIMAL(8,6)      NULL,
    max_lat       DECIMAL(8,6)      NULL,
    min_lon       DECIMAL(9,6)      NULL,
    max_lon       DECIMAL(9,6)      NULL,
    PRIMARY KEY (id),
    UNIQUE KEY uq_trip_code (trip_code),
    KEY idx_trip_taxi_start (taxi_id, start_time),
    KEY idx_trip_call_type (call_type),
    CONSTRAINT fk_trip_taxi FOREIGN KEY (taxi_id)
        REFERENCES taxi (taxi_id) ON DELETE CASCADE ON UPDATE CASCADE,
    CONSTRAINT fk_trip_call_type FOREIGN KEY (call_type)
        REFERENCES call_type (code) ON DELETE RESTRICT ON UPDATE CASCADE,
    CONSTRAINT chk_trip_day_type CHECK (day_type IN ('A', 'B', 'C')),
    CONSTRAINT chk_trip_times    CHECK (end_time >= start_time)
);

CREATE TABLE gps_point (
    trip_id  INT UNSIGNED      NOT NULL,
    seq      SMALLINT UNSIGNED NOT NULL,
    lat      DECIMAL(8,6)      NOT NULL,
    lon      DECIMAL(9,6)      NOT NULL,
    PRIMARY KEY (trip_id, seq),
    CONSTRAINT fk_gps_point_trip FOREIGN KEY (trip_id)
        REFERENCES trip (id) ON DELETE CASCADE ON UPDATE CASCADE
);
\end{lstlisting}

\subsection*{Part 2 -- Queries}

The queries are in two Python programs, \texttt{queries/part2\_q1\_5.py} and
\texttt{queries/part2\_q6\_10.py}, built like \texttt{example.py} from the
assignment. Each program opens a connection through \texttt{DbConnector}
(\texttt{mysql-connector-python}, settings from a \texttt{.env} file), runs the
SQL for each task with \texttt{cursor.execute()} and reads the rows with
\texttt{cursor.fetchall()}. All answers are computed in MySQL. Python only
prints the results as tables with \texttt{tabulate}, does a few small counts
on the returned rows (for example how many circular trips hardly move in
Q9), and writes the full lists to \texttt{queries/output/}. One task can be
run on its own, for example
\texttt{python queries/part2\_q6\_10.py 6}.

The two queries that do the most work are Q6 and Q10. Q6 first uses the
bounding box stored on each trip to skip trips that are nowhere near City
Hall, and then measures the distance from each remaining GPS point (inner
part of the query). The margins of $\pm$0.0010$^\circ$ latitude and
$\pm$0.0013$^\circ$ longitude are about 110\,m at Porto, so the box
contains the whole 100\,m circle:
\begin{lstlisting}[language=SQL]
SELECT p.trip_id,
       MIN(ST_Distance_Sphere(POINT(p.lon, p.lat),
                              POINT(-8.62911, 41.15794))) AS min_dist_m
FROM trip t
JOIN gps_point p ON p.trip_id = t.id
WHERE t.min_lat <= 41.15794 + 0.0010 AND t.max_lat >= 41.15794 - 0.0010
  AND t.min_lon <= -8.62911 + 0.0013 AND t.max_lon >= -8.62911 - 0.0013
  AND p.lat BETWEEN 41.15794 - 0.0010 AND 41.15794 + 0.0010
  AND p.lon BETWEEN -8.62911 - 0.0013 AND -8.62911 + 0.0013
GROUP BY p.trip_id
HAVING min_dist_m <= 100;
\end{lstlisting}
Q6 and Q9 use MySQL's \texttt{ST\_Distance\_Sphere} instead of the
\texttt{haversine} package, so the distances are computed in the database
and the GPS points never have to be sent to Python. Both compute the
great-circle distance on a sphere, and the Q6 answer matches the haversine
count from the EDA.

Q10 uses the window function \texttt{LAG} to find the end of the previous
trip of the same taxi. The idle time is the start of a trip minus the end of
the previous one. We compute it in Unix seconds rather than local
\texttt{DATETIME}, so idle times across the two daylight saving changes are
not off by an hour. The version below is shortened. The program also
computes the 8-hour variant and counts overlapping trips:
\begin{lstlisting}[language=SQL]
WITH ordered AS (
    SELECT taxi_id,
           CAST(start_unix AS SIGNED)
             - CAST(LAG(start_unix + duration_s)
                    OVER (PARTITION BY taxi_id ORDER BY start_unix, id)
               AS SIGNED) AS idle_s
    FROM trip
    WHERE is_valid = 1
)
SELECT taxi_id, ROUND(AVG(idle_s) / 3600, 2) AS avg_idle_h
FROM ordered
WHERE idle_s >= 0
GROUP BY taxi_id
ORDER BY avg_idle_h DESC
LIMIT 20;
\end{lstlisting}
The full SQL for every task is in the two programs.

The assignment defines an invalid trip as a trip with fewer than 3 GPS
points, so for the counting questions Q1--Q4a we count every trip in the
database. From Q4b on the questions are about durations, distances,
positions and the time between trips, where a trip with 0--2 points has no
usable value, so we only use valid trips (at least 3 points). Such a trip
lasts at most 15\,s and has no real route, so it would pull the averages in
Q4b and Q5 down, and a trip with one point would always count as circular
in Q9. All times are local Porto time.

\textbf{Q1.} There are \textbf{448 taxis, 1,710,557 trips and 83,408,413 GPS
points}. Counting only valid trips, there are 1,666,742 trips with
83,363,103 points.
\screenshot{q1}{Q1. The last column of the second table is the
\texttt{SUM(n\_points)} cross-check.}

\textbf{Q2.} On average each taxi has \textbf{3,818.2 trips} (3,720.4 if we
only count valid trips). The least active taxi has 1 trip, the most active
10,731.
\screenshot{q2}{Q2.}

\textbf{Q3.} Taxi 20000080 has the most trips (10,731), followed by 20000403
(9,236) and 20000066 (8,443). Number 20 is 20000235 with 6,410 trips.
61\,\% of the top taxi's trips have fewer than 3 points. With valid trips
only, 20000080 (4,211 valid trips) would not be in the top 20.
\screenshot{q3}{Q3: the 20 taxis with the most trips.}

\textbf{Q4a.} B (taxi stand) is the most used call type for 354 taxis, C for
82 and A for 12. Ties are broken alphabetically, so the one tied taxi
(20000344, as many B as C trips) gets B. The screenshot shows the first 20
taxis. The full table for all 448 taxis is in
\texttt{queries/output/q4a\_full.txt}.

\textbf{Q4b.} Street hails (C) are the longest trips on average, and they are
 more likely to start at night: 32\,\% of C trips start between 00 and
06, against 12--13\,\% for A and B (Table~\ref{tab:q4b}).

\begin{table}[H]\centering\small
\begin{tabular}{lrrrrrr}
\toprule
 & Avg.\ duration & Avg.\ distance & \multicolumn{4}{c}{Share of trips starting}\\
\cmidrule(lr){4-7}
Call type & (min) & (km) & 00--06 & 06--12 & 12--18 & 18--24\\
\midrule
A (central) & 12.59 & 5.43 & 12.05\,\% & 31.82\,\% & 32.59\,\% & 23.55\,\%\\
B (stand)   & 11.23 & 5.11 & 13.02\,\% & 27.76\,\% & 34.56\,\% & 24.65\,\%\\
C (street)  & 13.67 & 6.67 & 31.78\,\% & 24.40\,\% & 24.00\,\% & 19.82\,\%\\
\bottomrule
\end{tabular}
\caption{Q4b: average duration and distance, and share of trips starting in
each time band, per call type (valid trips, local time).}
\label{tab:q4b}
\end{table}
Without the trips that have a GPS jump, the average distances are 1.5--4\,\%
lower (A 5.34, B 5.01, C 6.40\,km).
\screenshot{q4ab}{Q4a (first 20 taxis and summary) and Q4b.}

\textbf{Q5.} We define hours driven as the sum of the trip durations, and
distance driven as the sum of \texttt{distance\_km} (haversine between
consecutive points, computed when loading). Taxi \textbf{20000904} has both
the most hours (1,961.0\,h) and the longest total distance (62,788\,km).
Second on hours is 20000129
(1,649.6\,h). Second on distance is 20000351 (56,639\,km), which is only
number 51 on hours. Without the trips that have a GPS jump, 20000904 is
still first on both, and 20000351 is still second on distance, so the answer
does not depend on these trips.
\screenshot{q5}{Q5.}

\textbf{Q6.} \textbf{126,295} valid trips pass within 100\,m of Porto City
Hall. This matches what we found directly from the CSV during the EDA. The
screenshot shows the first 20. The full list is in
\texttt{queries/output/q6\_trips.csv}. It has 126,532 rows, since it also
includes the 237 invalid trips (marked with \texttt{is\_valid = 0}).
\screenshot{q6}{Q6.}

\textbf{Q7.} There are \textbf{43,815 invalid trips} (2.56\,\%): 5,894 with
no points, 30,532 with one and 7,389 with two.
\screenshot{q7}{Q7.}

\textbf{Q8.} A trip is a midnight crosser if
\texttt{DATE(end\_time) > DATE(start\_time)} in local time.
\textbf{8,028} trips start on one day and end on the next (0.48\,\% of valid
trips). None of them crosses more than one midnight. The full list is in
\texttt{queries/output/q8\_midnight\_trips.csv}.
\screenshot{q8}{Q8.}

\textbf{Q9.} \textbf{21,467} trips start and end within 50\,m of each other,
measured between the stored start and end point. 7,608 of them hardly move at all (less than 100\,m driven), while 12,177
drive at least 1\,km before coming back. We still count the trips that
hardly move, since the task only asks for start and end within 50\,m. The
full list is in \texttt{queries/output/q9\_circular\_trips.csv}.
\screenshot{q9}{Q9.}

\textbf{Q10.} The two taxis at the top, 20000941 (1,746.6\,h) and 20000969
(282.1\,h), had almost no trips (4 and 25 valid trips). From number 3
(20000312, 13.1\,h) the averages are between 5 and 13 hours. As in the other
queries from Q4b on, we only use valid trips. Most invalid trips have 0 or
1 point, and some taxis have thousands of them (taxi 20000080 has 6,520), so
they look like recording errors, and including them would split real breaks
into many short ones. In 792 cases a trip starts before the previous trip of
the same taxi has ended. A negative idle time is not a break, so we leave
these cases out of the average. They are only 0.05\,\% of all gaps.
\screenshot{q10}{Q10: the 20 taxis with the highest average idle time.}

% =====================================================================
\section*{Discussion}

\textbf{What we did differently.} Group 16 has access to a VM at IDI, but we
ran MySQL 8.0.39 locally in Docker instead. It was simpler to set up, and
our own computer had enough power to load and query the 83 million GPS
points.
The assignment does not say which timezone to use. We chose local Porto
time, because that is when the trips happened. 157,929 trips (9\,\%) end up
in a different time band than with UTC, but over all trips the share in
each band changes by at most 1.1 percentage points.
We also did not use TRIP\_ID as the primary key, even though the assignment
calls it a unique identifier: it is not unique in the raw data, and a
smaller key of our own saves space (see Part 1 -- Schema).

\textbf{Invalid trips and GPS jumps.} We kept both and marked them instead
of deleting them. This way Q7 can count the invalid trips, and the other
queries can choose what to include. The downside is that every query has to
state which trips it uses. For distances the difference turned out to be
small, only 1.5--4\,\% in Q4b.

\textbf{The flags in the dataset are not reliable.} DAY\_TYPE is always ``A''
and MISSING\_DATA marks only 10 trips, while about 38,000 trips clearly have
gaps. We therefore judged data quality from the trajectories themselves.

\textbf{Rankings per taxi.} Taxis with few or strange trips end up at the top
of some rankings. The top taxi in Q3 has 61\,\% invalid trips, and the two
highest values in Q10 come from taxis with only 4 and 25 trips. In Q10 we
define idle time as the next start minus the previous end. Without a limit,
days off dominate the average. If we only count breaks shorter than 8 hours,
the top 20 are between 85 and 106 minutes. We still use the plain average
as our main answer, since that is what the task asks for, and show the
8-hour variant next to it in the output.

\textbf{Limitations.} Q6 only checks the GPS points, not the line between
them. At 50\,km/h a taxi moves about 210\,m in 15\,s, which is more than the
200\,m width of the circle, so a taxi can pass City Hall without any point
inside it. We also compute the duration as 15\,s times the number of steps
between points. If a jump is a gap in the signal rather than a bad point,
the trip actually took longer, so its duration and end time are too short.
This affects the 38,198 trips with a jump (2.2\,\%). We cannot correct these
durations, but leaving the trips out changes the average durations in Q4b
by at most 0.2 minutes and does not change the top taxi in Q5, so we do not
expect the error to change our conclusions.

\textbf{Pain points.} The main practical problem was the size of the data.
To keep the load time and disk use down with 83 million GPS points, we
bulk-loaded with \texttt{LOAD DATA}, turned off the binary log and the
foreign key checks during the load, and kept \texttt{gps\_point} as narrow
as possible.

\textbf{What we learned.} We learned that it pays off to explore the data
properly before designing the schema. Both the duplicate IDs and the GPS
jumps would have caused problems if we had found them later. We also learned
how much schema choices such as key size and precomputed columns matter
when a table has tens of millions of rows.

% =====================================================================
\section*{Feedback}

\begin{itemize}
  \item The dataset description does not match the data: TRIP\_ID is not
        unique, DAYTYPE is ``A'' in every row, and MISSING\_DATA marks only
        10 trips.
  \item The time bands in Q4b and the midnight in Q8 depend on the
        timezone, which the assignment does not specify.
\end{itemize}

% =====================================================================
\section*{AI use}

We used Claude Code, an AI coding assistant, throughout the assignment.

\textbf{Setup and planning.} We created a project in Claude Code and first
sat down together to collect the sources the AI should work from: the
assignment text, the report template and NTNU's guidelines on AI in student
work. This gave the AI the same starting point as us and helped us keep
control of the project. Based on these sources we used the AI to sketch out
what we needed: the EDA checks, the cleaning rules, the schema, the loader and
the ten queries. We then worked through that plan step by step.

\textbf{Code.} Working together, we had the AI write the first versions of
the EDA script, the schema, the loader and the query programs from our
instructions. That first version
was a good start, but far from the final result. Both of us went back and
forth with the AI, asked why things were done the way they were, and
questioned several of the choices, for example how to handle duplicate trip
IDs, whether to delete or keep the short trips, and which timezone to use.
It took many rounds of critical review to reach the final version, and
because of that we understand and can explain all the code we hand in, and
can justify every design choice.

\textbf{Checking the results.} We ran all the code ourselves and checked the
answers by comparing the query results with the EDA and by looking at the
data directly in the database.

\textbf{Report.} The AI helped us draft, structure and review the report
text, and computed some of the supporting numbers from the data. We went
through all of it, rewrote parts of it and checked every number against the
output.

\end{document}
